\documentclass[11pt]{article}
\usepackage[a4paper,margin=21mm]{geometry}
\usepackage[no-math]{fontspec}
\setmainfont{Latin Modern Roman}
\newfontfamily\CJKfont{Noto Serif CJK SC}
\XeTeXlinebreaklocale "zh"
\XeTeXlinebreakskip=0pt plus 1pt
\usepackage{amsmath,amssymb,amsthm,amscd,graphicx,color}
\usepackage{xurl}
\usepackage[unicode,hidelinks]{hyperref}
\allowdisplaybreaks
\setlength\emergencystretch{3em}
\numberwithin{equation}{section}

\usepackage{mathrsfs}
\usepackage{tikz-cd}

\newtheorem{Theorem}{Theorem}[section]
\newtheorem*{Theorem*}{Theorem}
\newtheorem{Corollary}[Theorem]{Corollary}
\newtheorem{Lemma}[Theorem]{Lemma}
\newtheorem{Proposition}[Theorem]{Proposition}
{ \theoremstyle{definition}
\newtheorem{Definition}[Theorem]{Definition}
\newtheorem{Note}[Theorem]{Note}
\newtheorem{Construction}[Theorem]{Construction}
\newtheorem{Example}[Theorem]{Example}
\newtheorem{Remark}[Theorem]{Remark} }

\numberwithin{equation}{section}
\newfontfamily\nativebbk{DejaVu Math TeX Gyre}
\ExplSyntaxOn
\NewDocumentCommand{\mathds}{m}{\str_if_eq:nnTF{#1}{k}{\mathord{\text{\mdseries\upshape\nativebbk\char"1D55C}}}{\PackageError{nativecompat}{Unsupported~mathds~argument}{Only~verified~lowercase~k~is~accepted}}}
\ExplSyntaxOff
\setmainfont{Noto Serif}

\begin{document}
\begin{center}\Large Homotopy classification of faithful finite-group graded extensions of fusion 2-categories over an algebraically closed field of characteristic zero\end{center}
\noindent\textbf{Stable ID:} 1050\quad\textbf{Author:} Thibault Didier Décoppet\par
\noindent\textbf{Work:} Extension Theory and Fermionic Strongly Fusion 2-Categories\par
\noindent\textbf{Source:} \url{https://sigma-journal.com/2024/092/}\par
\noindent\textbf{Version:} Official SIGMA 20 (2024) article 092, DOI 10.3842/SIGMA.2024.092; journal PDF and native TeX acquired 2026-09-30, exact bytes pinned by SHA256\par
\noindent\textbf{Rights:} CC BY 4.0. Authors retain copyright. \url{https://creativecommons.org/licenses/by/4.0/}. Reuse and typesetting adaptation; original language and mathematical content preserved.\par
\noindent\textbf{Difficulty (prior selection preserved):} {\CJKfont H2: The proof converts a homotopy map into a strongly monoidal functor and constructs the graded algebra from the coproduct completion of the group. The key additional obligation is to establish rigidity of the resulting monoidal 2-category by representing adjoint module functors in inverse grading components. The reverse direction requires the full relative tensor-product equivalence of graded components, supplied by the bundled Proposition3.4.}\par
\medskip\noindent\textbf{Editorial boundary note (not author text):} Complete original Theorem3.11 over an algebraically closed characteristic-zero field, classifying faithful finite-group extensions by homotopy maps to the delooped Brauer–Picard space. Full current fusion/module/relative tensor and invertibility definitions, graded-component tensor equivalence Proposition3.4 with its entire human bimodule-algebra proof, inverse-component duality and Brauer–Picard definition are included. The complete proof retains the coproduct-completion/lax-functor factorization (native editable tikz-cd), graded monoidal construction, representation of adjoint module functors in inverse grading components establishing rigidity, analogous right-dual author wording and the full reverse strongly monoidal/invertible-components direction. The exact pre-existing Morita-dual and relative tensor-product theorems and formal monoidal/coherence prerequisites remain precisely author-referenced. One primary extension classification, positioned as the first main result; fermionic applications, examples and supercohomology are excluded and not counted. No assistant proof, missing-coherence reconstruction or mathematical review. The unavailable dsfont import is replaced only by an exact lowercase-k-only formatting resolver using the installed named DejaVu Math TeX Gyre U+1D55C double-struck field glyph. Every literal author mathds\{k\} call remains unchanged, every active argument is verified k, and all other arguments cause an error. This named font is a required portable dependency; no ordinary k or uppercase K is substituted. Wrapper prose uses installed named Noto Serif so the literal bibliography cedilla in Pstr\textbackslash{}c\{a\}gowski is rendered rather than dropped; the author accent command is not changed or converted to an ogonek. Standard original-style math fonts remain separate under fontspec no-math.\par
\medskip\hrule\medskip\noindent\textbf{Author text begins below.}\par
\setcounter{section}{1}
% Original source lines 86--191
\section{Preliminaries}

\subsection{Fusion 2-categories}

Over an algebraically closed field $\mathds{k}$ of characteristic zero, the notions of finite semisimple 2-category and fusion 2-category were introduced in \cite{DR}. We succinctly review these definitions here, and refer the reader to the aforementioned reference and also \cite{D2} for details.

A ($\mathds{k}$-linear) 2-category $\mathfrak{C}$ is locally finite semisimple if its $\operatorname{Hom}$-1-categories are finite semisimple. An object $C$ of $\mathfrak{C}$ is called simple if $\operatorname{Id}_C$ is a simple object of the finite semisimple 1-category~$\operatorname{End}_{\mathfrak{C}}(C)$. Then, a \textit{finite semisimple} 2-category is a locally finite semisimple 2-category that is Cauchy complete in the sense of \cite{GJF} (see also \cite{D1}), has right and left adjoints for 1-morphisms, and has finitely many equivalence classes of simple objects. In a finite semisimple 2-category $\mathfrak{C}$, any two simple objects are called connected if there exists a non-zero 1-morphism between them. This defines an equivalence relation, and we write $\pi_0(\mathfrak{C})$ for the corresponding set of equivalences classes; This is the set of \textit{connected components} of $\mathfrak{C}$.

Now, the definition of a monoidal 2-category is well known. In particular, using the notations of \cite{SP}, a monoidal structure on the 2-category $\mathfrak{C}$ involves a 2-functor $\Box\colon \mathfrak{C}\times \mathfrak{C}\rightarrow\mathfrak{C}$, the monoidal product, and a distinguished object $I$, the monoidal unit, together with various other coherence data. Given any object $C$ of $\mathfrak{C}$, we can ask for the existence of its left dual $^{\sharp}C$ and its right dual~$C^{\sharp}$. These two notions were studied in \cite{Pstr}. Categorifying the concept of a fusion 1-category is therefore straightforward, and we obtain the following definition.

\begin{Definition}
A multifusion 2-category is a finite semisimple monoidal 2-category $\mathfrak{C}$ such that every object $C$ admits a right dual $C^{\sharp}$ and a left dual $^{\sharp}C$. A fusion 2-category is a multifusion 2-category whose monoidal unit $I$ is a simple object.
\end{Definition}

To every fusion 2-category $\mathfrak{C}$, there is an associated braided fusion 1-category $\Omega\mathfrak{C}:=\operatorname{End}_{\mathfrak{C}}(I)$, the 1-category of endomorphisms of the monoidal unit $I$ of $\mathfrak{C}$. Conversely, given any braided fusion 1-category $\mathcal{B}$, we can consider the fusion 2-category $\mathbf{Mod}(\mathcal{B})$ of finite semisimple left $\mathcal{B}$-module 1-categories with monoidal structure given by $\boxtimes_{\mathcal{B}}$, the relative Deligne tensor product. Below, we give another class of examples that will be relevant for our purposes. Many more may be found in \cite{DR} and \cite{DY23}, and others will be discussed subsequently.

\begin{Example}
Given $G$ a finite group, and $\pi$ a 4-cocycle for $G$ with coefficients in $\mathds{k}^{\times}$, we can then consider the fusion 2-category $\mathbf{2Vect}^{\pi}(G)$ of $G$-graded finite semisimple 1-categories, also known as 2-vector spaces, with pentagonator twisted by $\pi$. Such fusion 2-category were completely characterized in \cite{JFY} as those fusion 2-categories $\mathfrak{C}$ for which $\Omega\mathfrak{C}\simeq \mathbf{Vect}$. We call these bosonic strongly fusion 2-categories.
\end{Example}

\subsection{Module 2-categories and the relative 2-Deligne tensor product}

Let $\mathfrak{C}$ be a monoidal 2-category with monoidal product $\Box$, and monoidal unit $I$. A right $\mathfrak{C}$-module 2-category is a 2-category $\mathfrak{M}$ equipped with an action 2-functor $\Box\colon\mathfrak{M}\times\mathfrak{C}\rightarrow \mathfrak{M}$, and coherence data satisfying various axioms. Likewise, there is a notion of left $\mathfrak{C}$-module 2-category. Additionally, there are notions of $\mathfrak{C}$-module 2-functors, $\mathfrak{C}$-module 2-natural transformations, and~$\mathfrak{C}$-module modifications. There are also notions of bimodule 2-categories, and maps between them. We refer the reader to \cite[Section~2]{D4} for the precise definitions.

Let us now assume that $\mathfrak{C}$ is a fusion 2-category. If $\mathfrak{M}$ is a left $\mathfrak{C}$-module 2-category, then we can consider $\mathbf{\operatorname{End}}_{\mathfrak{C}}(\mathfrak{M})$, the monoidal 2-category of left $\mathfrak{C}$-module 2-endofunctors on $\mathfrak{M}$. The next result combines \cite[Theorem~5.3.2]{D8} with \cite[Corollary~5.1.3]{D9}.

\begin{Theorem}\label{thm:Moritadual}
The monoidal $2$-category $\mathbf{\operatorname{End}}_{\mathfrak{C}}(\mathfrak{M})$ is a multifusion $2$-category.
\end{Theorem}

Moreover, the multifusion 2-category $\mathbf{\operatorname{End}}_{\mathfrak{C}}(\mathfrak{M})$ is fusion if and only if $\mathfrak{M}$ is indecomposable as a finite semisimple left $\mathfrak{C}$-module 2-category.

For our purposes, we need to consider another operation on finite semisimple module 2-categories. Let $\mathfrak{M}$ be a finite semisimple right $\mathfrak{C}$-module 2-category and $\mathfrak{N}$ be a finite semisimple left $\mathfrak{C}$-module 2-category. There are notions of $\mathfrak{C}$-balanced 2-functors, $\mathfrak{C}$-balanced 2-natural transformations, and $\mathfrak{C}$-balanced modifications out of $\mathfrak{M}\times\mathfrak{N}$. We refer the reader to \cite[Section~2.2]{D10} for the details. The next result is a combination of \cite[Theorem~2.2.4]{D10} with \cite[Corollary~5.1.3]{D9}.

\begin{Theorem}
There is a $3$-universal $\mathfrak{C}$-balanced $2$-functor $\boxtimes_{\mathfrak{C}}\colon \mathfrak{M}\times\mathfrak{N}\rightarrow \mathfrak{M}\boxtimes_{\mathfrak{C}}\mathfrak{N}$ to a finite semisimple $2$-category.
\end{Theorem}

In fact, the proof of \cite[Theorem~2.2.4]{D10} gives an explicit description of $\mathfrak{M}\boxtimes_{\mathfrak{C}}\mathfrak{N}$. More precisely, if $A$ is a (necessarily separable) algebra in $\mathfrak{C}$ such that $\mathfrak{M}$ is equivalent to $\mathbf{LMod}_{\mathfrak{C}}(A)$, the right~$\mathfrak{C}$-module 2-category of left $A$-modules in $\mathfrak{C}$, and $B$ is a (necessarily separable) algebra in $\mathfrak{C}$ such that $\mathfrak{N}$ is equivalent to $\mathbf{Mod}_{\mathfrak{C}}(B)$, the left $\mathfrak{C}$-module 2-category of right $B$-modules in~$\mathfrak{C}$, then~$\mathfrak{M}\boxtimes_{\mathfrak{C}}\mathfrak{N}\simeq \mathbf{Bimod}_{\mathfrak{C}}(A,B)$, the finite semisimple 2-category of $A$-$B$-bimodules in $\mathfrak{C}$.

Finally, let us recall that, as was explained in \cite[Section~3.2]{D10}, the existence of the relative 2-Deligne tensor product allows us to consider the symmetric monoidal Morita 4-category $\mathbf{F2C}$ of (multi)fusion 2-categories, finite semisimple bimodule 2-categories, and bimodule morphisms.

\section{Extension theory}

For fusion 1-categories over an algebraically closed field of characteristic zero, extension theory was developed in \cite{ENO2}. The key concept is that of an invertible bimodule 1-category, or, more precisely, the space formed by such objects together with their invertible morphisms. Proceeding in a similar fashion, we will discuss extension theory for fusion 2-categories, i.e., we will study group graded fusion 2-categories in the sense of the definition below.

\begin{Definition}
Let $\mathfrak{C}$ be a fusion 2-category and $G$ a finite group. A $G$-grading on $\mathfrak{C}$ is a~decomposition $\boxplus_{g\in G}\mathfrak{C}_g$ of $\mathfrak{C}$ into a direct sum of finite semisimple 2-categories such that for every $C$ in $\mathfrak{C}_g$ and $D$ in $\mathfrak{C}_h$, $C\Box D$ lies $\mathfrak{C}_{gh}$. A $G$-grading on $\mathfrak{C}$ is faithful if $\mathfrak{C}_g$ is non-zero for every $g$ in $G$.
\end{Definition}

\subsection{Invertible bimodule 2-categories}

Let $\mathds{k}$ be an algebraically closed field of characteristic zero. We begin by recalling a definition from \cite{DY23}. Let $\mathfrak{C}$ and $\mathfrak{D}$ be two fusion 2-categories. We write $\mathfrak{D}^{\rm mop}$ for the fusion 2-category obtained by endowing the finite semisimple 2-category $\mathfrak{D}$ with the opposite of the monoidal product of $\mathfrak{D}$.

\begin{Definition}
A finite semisimple $\mathfrak{C}$-$\mathfrak{D}$-bimodule 2-category $\mathfrak{M}$ is invertible if the canonical monoidal 2-functor $\mathfrak{D}^{\rm mop}\rightarrow \mathbf{\operatorname{End}}_{\mathfrak{C}}(\mathfrak{M})$ is an equivalence.
\end{Definition}

At the decategorified level, i.e., for fusion 1-categories, invertibility of a finite semisimple bimodule 1-category admits various equivalent characterizations as is explained in \cite[Proposition~4.2]{ENO2}. A categorified version of this result was given in \cite{D10}, which we partially recall below.

\begin{Proposition}\label{prop:invertibilitycharacterization}
Let $\mathfrak{M}$ be a finite semisimple $\mathfrak{C}$-$\mathfrak{D}$-bimodule $2$-category. The following are equivalent:
\begin{enumerate}\itemsep=0pt
 \item[$(1)$] The $\mathfrak{C}$-$\mathfrak{D}$-bimodule $2$-category $\mathfrak{M}$ is invertible.
 \item[$(2)$] The $\mathfrak{C}$-$\mathfrak{D}$-bimodule $2$-category $\mathfrak{M}$ defines an invertible $1$-morphism from $\mathfrak{C}$ to $\mathfrak{D}$ in $\mathbf{F2C}$.
\end{enumerate}
\end{Proposition}

The next proposition gives preliminary insight into the relation between group graded fusion 2-categories and invertible bimodule 2-categories. We note that the first part of the result below has already appeared as \cite[Proposition~3.1.7]{DY23}. We also refer the reader to \cite[Theorem~6.1]{ENO2} for the decategorified version of this proposition.

\begin{Proposition}\label{prop:gradedMoritaequivalence}
Let $\mathfrak{C}$ be a faithfully $G$-graded fusion $2$-category $\mathfrak{C}$. For any $g\in G$, the finite semisimple $\mathfrak{C}_e$-$\mathfrak{C}_e$-bimodule $2$-category $\mathfrak{C}_g$ is invertible. Furthermore, for any $g,h\in G$, the $2$-functor $\Box\colon\mathfrak{C}\times\mathfrak{C}\rightarrow \mathfrak{C}$ induces an equivalence \smash{$\mathfrak{C}_g\boxtimes_{\mathfrak{C}_e}\mathfrak{C}_h\xrightarrow{\simeq}\mathfrak{C}_{gh}$} of $\mathfrak{C}_e$-$\mathfrak{C}_e$-bimodule $2$-categories.
\end{Proposition}
\begin{proof}
The first part follows from the second, so we only prove the latter. For the second part, note that it is enough to show that the canonical 2-functor is an equivalence. To this end, pick any non-zero objects $X$ in $\mathfrak{C}_{g^{-1}}$, and $Y$ in $\mathfrak{C}_{h}$. It follows from the proof of \cite[Theorem~5.4.3]{D8} that~$Y\Box {^{\sharp}Y}$ is a separable algebra in $\mathfrak{C}$, and that the left $\mathfrak{C}$-module 2-functor $\mathfrak{C}\rightarrow \mathbf{Mod}_{\mathfrak{C}}\big(Y\Box {^{\sharp}Y}\big)$ given by $C\mapsto C\Box {^{\sharp}Y}$ is an equivalence. Likewise, $X\Box {^{\sharp}X}$ is a separable algebra in $\mathfrak{C}$, and the right~$\mathfrak{C}$-module 2-functor $\mathfrak{C}\rightarrow \mathbf{LMod}_{\mathfrak{C}}\big(X\Box {^{\sharp}X}\big)$ given by $C\mapsto X\Box C$ is an equivalence. In particular, we also find that the 2-functor $\mathfrak{C}\rightarrow \mathbf{Bimod}_{\mathfrak{C}}\big(X\Box {^{\sharp}X},Y\Box {^{\sharp}Y}\big)$ given by $C\mapsto X\Box C\Box {^{\sharp}Y}$ is an equivalence. Under these identifications, as was recalled above, it follows from the proof of~\cite[Theorem~2.2.4]{D10} that the canonical 2-functor $\mathfrak{C}\times\mathfrak{C}\rightarrow\mathfrak{C}\boxtimes_{\mathfrak{C}}\mathfrak{C}$ is identified with $\Box\colon\mathfrak{C}\times\mathfrak{C}\rightarrow \mathfrak{C}$, which is given by
\begin{gather*}
 \mathbf{LMod}_{\mathfrak{C}}\big(X\Box {^{\sharp}X}\big)\times\mathbf{Mod}_{\mathfrak{C}}\big(Y\Box {^{\sharp}Y}\big) \rightarrow \mathbf{Bimod}_{\mathfrak{C}}\big(X\Box {^{\sharp}X},Y\Box {^{\sharp}Y}\big),\qquad
(M,N) \mapsto M\Box N.
\end{gather*}
 But, by considering the restriction $\mathfrak{C}_e\hookrightarrow\mathfrak{C}$, the above equivalences restrict to equivalences $\mathfrak{C}_h\rightarrow \mathbf{Mod}_{\mathfrak{C}_e}\big(Y\Box {^{\sharp}Y}\big)$, $\mathfrak{C}_g\rightarrow \mathbf{LMod}_{\mathfrak{C}_e}\big(X\Box {^{\sharp}X}\big)$, and $\mathfrak{C}_{gh}\rightarrow \mathbf{Bimod}_{\mathfrak{C}_e}\big(X\Box {^{\sharp}X},Y\Box {^{\sharp}Y}\big)$. In particular, the canonical 2-functor $\mathfrak{C}_g\times\mathfrak{C}_h\rightarrow\mathfrak{C}_g\boxtimes_{\mathfrak{C}_e}\mathfrak{C}_h$ is identified with $\Box:\mathfrak{C}_g\times\mathfrak{C}_h\rightarrow\mathfrak{C}_{gh}$ as claimed.
\end{proof}

A similar argument as the one used in the proof of the above proposition yields the following result.

\begin{Corollary}\label{cor:invertibledual}
The canonical $\mathfrak{C}_e$-$\mathfrak{C}_e$-bimodule $2$-functor $\mathfrak{C}_{g^{-1}}\rightarrow \mathbf{Fun}_{\mathfrak{C}_e}(\mathfrak{C}_g,\mathfrak{C}_e)$ given by $D\mapsto \{C\mapsto D\Box C\}$ is an equivalence.
\end{Corollary}

\subsection{Brauer--Picard spaces and extensions}

Group-graded extensions of fusion 1-categories are parameterised by maps into the space of invertible bimodule 1-categories and their (invertible) higher morphisms \cite{ENO2}. Thanks to the existence of the 4-category $\mathbf{F2C}$ obtained in \cite{D10}, the corresponding spaces for our 2-categorical purposes are easy to define.

\begin{Definition}
Let $\mathfrak{C}$ be a fusion 2-category. The Brauer--Picard space of $\mathfrak{C}$ consists of the invertible objects and the invertible morphisms in the monoidal 3-category of finite semisimple $\mathfrak{C}$-$\mathfrak{C}$-bimodule 2-categories, that is, \begin{gather*}\mathscr{B}{\rm r}\mathscr{P}{\rm ic}(\mathfrak{C}):=\big(\mathbf{\operatorname{End}}_{\mathbf{F2C}}(\mathfrak{C})\big)^{\times}.\end{gather*}
\end{Definition}

\begin{Remark}
We will write $Br\operatorname{Pic}(\mathfrak{C}):= \pi_0(\mathscr{B}{\rm r}\mathscr{P}{\rm ic}(\mathfrak{C}))$, the Brauer--Picard group of $\mathfrak{C}$. It follows from \cite[Lemma~2.2.1]{D9} that $\Omega \mathbf{\operatorname{End}}_{\mathbf{F2C}}(\mathfrak{C})\simeq \mathscr{Z}(\mathfrak{C})$, the Drinfeld center of $\mathfrak{C}$, as defined in~\cite{Cr}. Thus, the homotopy groups of the space $\mathscr{B}{\rm r}\mathscr{P}{\rm ic}(\mathfrak{C})$ are given as follows:
 $$\renewcommand{\arraystretch}{1.2} \begin{tabular}{|c|c|c|c|c|}
\hline
$\pi_0$ & $\pi_1$ & $\pi_2$ & $\pi_3$\\
\hline
${\rm Br}\operatorname{Pic}(\mathfrak{C})$ & $\operatorname{Inv}(\mathscr{Z}(\mathfrak{C}))$ & $\operatorname{Inv}(\Omega\mathscr{Z}(\mathfrak{C}))$ & $\mathds{k}^{\times}$\\
\hline
\end{tabular}$$
\end{Remark}

\setcounter{Theorem}{10}
% Original source lines 231--258
A more general version of the next result has been announced \cite{JFR3}. We are very grateful to them for outlining their proof to us, which has inspired the argument that we give below.

\begin{Theorem}\label{thm:extensiontheory}
Let $G$ be a finite group, and $\mathfrak{C}$ be a fusion $2$-category. Then, faithfully $G$-graded extensions of $\mathfrak{C}$ are classified by homotopy classes of maps $\mathrm{B}G\rightarrow \mathrm{B}\mathscr{B}{\rm r}\mathscr{P}{\rm ic}(\mathfrak{C})$.
\end{Theorem}
\begin{proof}
For the purpose of this proof, it will be convenient to think of maps of spaces $\mathrm{B}G\rightarrow \mathrm{B}\mathscr{B}{\rm r}\mathscr{P}{\rm ic}(\mathfrak{C})$ as monoidal maps $G\rightarrow \mathscr{B}{\rm r}\mathscr{P}{\rm ic}(\mathfrak{C})$ of $(\infty,1)$-categories. We will also consider the $(\infty,1)$-category $\mathscr{C}$ obtained from $\mathbf{\operatorname{End}}_{\mathbf{F2C}}(\mathfrak{C})$ by only considering the invertible $n$-morphisms when $n\geq 2$. Then, by definition we have $\mathscr{B}{\rm r}\mathscr{P}{\rm ic}(\mathfrak{C})=\mathscr{C}^{\times}$ as monoidal $(\infty,1)$-categories.

We begin the proof by some general nonsense. Recall that algebras in the monoidal $(\infty,1)$-category $\mathscr{C}$ correspond precisely to lax monoidal functors $*\rightarrow \mathscr{C}$ (see \cite[Example~2.2.6.10]{L2}). We want a similar description for the notion of a (faithfully) $G$-graded algebra. In order to do so, consider the monoidal 1-category $G^{\sqcup}$, which is the coproduct completion of $G$. Now, there is a canonical algebra $A[G]$ in $G^{\sqcup}$ given by
 \begin{gather*}
A[G]:=\coprod_{g\in G}g,
\end{gather*}
 or equivalently a lax monoidal functor $*\rightarrow G^{\sqcup}$. Then, giving a $G$-grading on an algebra $A$ in $\mathscr{C}$ is equivalent to providing a factorization of lax monoidal functors
\begin{displaymath}
\begin{tikzcd}
* \arrow[rr,"A"] \arrow[rd, "A\lbrack G \rbrack"'] & & \mathscr{C} \\
 & G^{\sqcup}. \arrow[ru] &
\end{tikzcd}\end{displaymath}
But, given that $\mathscr{C}$ has direct sums, lax monoidal functors $G^{\sqcup}\rightarrow \mathscr{C}$ correspond exactly to lax monoidal functors $G\rightarrow \mathscr{C}$.

Now, given a map of spaces $\mathrm{B}G\rightarrow \mathrm{B}\mathscr{B}{\rm r}\mathscr{P}{\rm ic}(\mathfrak{C})$, or equivalently a strongly monoidal functor $F\colon G\rightarrow \mathscr{C}$, we obtain a $G$-graded algebra $\mathfrak{D}:=F(A[G])$ in $\mathfrak{C}$. We also write
\begin{gather*}
\mathfrak{D}=\boxplus_{g\in G}\mathfrak{D}_g=\boxplus_{g\in G}F(g).
\end{gather*}
 In fact, as $F(e) = \mathfrak{C}$ by definition, the monoidal 2-category $\mathfrak{D}$ is a faithfully $G$-graded extension of $\mathfrak{C}$. It is therefore enough to prove that $\mathfrak{D}$ is a fusion 2-category. The only property that is not obvious is rigidity. We have that $\mathbf{\operatorname{End}}_{\mathfrak{C}}(\mathfrak{D})$ is a fusion 2-category thanks to Theorem~\ref{thm:Moritadual} above. Let us fix $g\in G$. For any simple object $X$ in $\mathfrak{D}_g$, we can consider the left $\mathfrak{C}$-module 2-functor~${R_X\colon \mathfrak{D}\rightarrow \mathfrak{D}}$ given by $D\mapsto D\Box X$ on ${\mathfrak{D}_e\simeq \mathfrak{C}}$, and zero on $\mathfrak{D}_h$ for any~${e\neq h\in G}$. As the objects of the monoidal 2-category $\mathbf{\operatorname{End}}_{\mathfrak{C}}(\mathfrak{D})$ have duals, we can consider the right adjoint~$R_X^*$ of $R_X$. But, we have that $\mathfrak{D}_{g^{-1}}\rightarrow \mathbf{Fun}_{\mathfrak{D}_e}(\mathfrak{D}_g,\mathfrak{D}_e)$ is an equivalence of $\mathfrak{D}_e$-$\mathfrak{D}_e$-bimodule 2-categories as~$F$ is strongly monoidal, and therefore preserves duals. Thus, we find that ${R_X^*\simeq Y\Box (-)\colon \mathfrak{D}_g\rightarrow \mathfrak{D}_e=\mathfrak{C}}$ for some $Y$ in $\mathfrak{D}_{g^{-1}}$. This proves that $X$ has a left dual $Y$ in~$\mathfrak{D}$. One shows analogously that $X$ has a right dual.

Conversely, given $\mathfrak{D}=\boxplus_{g\in G}\mathfrak{D}_g$ a $G$-graded extension of $\mathfrak{C}= \mathfrak{D}_e$, we can consider the corresponding lax monoidal functor $F\colon G\rightarrow \mathscr{C}$. More precisely, we set $F(g) = \mathfrak{D}_g$ with lax monoidal structure given by the monoidal structure of $\mathfrak{D}$. We want to show that $F$ is strongly monoidal and factors through $\mathscr{B}{\rm r}\mathscr{P}{\rm ic}(\mathfrak{C})\subset \mathscr{C}$. It follows from the definition of an extension that $F$ is strongly unital. Proposition \ref{prop:gradedMoritaequivalence} above establishes that $F$ has the remaining desired properties.
\end{proof}
\begin{thebibliography}{99}
\bibitem[1]{Cr}
Crans S.E., Generalized centers of braided and sylleptic monoidal
 {$2$}-categories, \href{https://doi.org/10.1006/aima.1998.1720}{\textit{Adv. Math.}} \textbf{136} (1998), 183--223.

\bibitem[4]{D1}
D\'ecoppet T.D., Multifusion categories and finite semisimple 2-categories,
 \href{https://doi.org/10.1016/j.jpaa.2022.107029}{\textit{J.~Pure Appl. Algebra}} \textbf{226} (2022), 107029, 16~pages,
 \href{https://arxiv.org/abs/2012.15774}{arXiv:2012.15774}.

\bibitem[5]{D2}
D\'ecoppet T.D., Weak fusion 2-categories, \textit{Cah. Topol. G\'eom.
 Diff\'er. Cat\'eg.} \textbf{63} (2022), 3--24, \href{https://arxiv.org/abs/2103.15150}{arXiv:2103.15150}.

\bibitem[6]{D8}
D\'ecoppet T.D., The {M}orita theory of fusion 2-categories, \textit{High.
 Struct.} \textbf{7} (2023), 234--292, \href{https://arxiv.org/abs/2208.08722}{arXiv:2208.08722}.

\bibitem[7]{D9}
D\'ecoppet T.D., Drinfeld centers and {M}orita equivalence classes of fusion
 2-categories, \textit{Compos. Math.}, {t}o appear, \href{https://arxiv.org/abs/2211.04917}{arXiv:2211.04917}.

\bibitem[8]{D4}
D\'ecoppet T.D., Finite semisimple module 2-categories, \textit{Selecta Math.~(N.S.)}, {t}o appear, \href{https://arxiv.org/abs/2107.11037}{arXiv:2107.11037}.

\bibitem[9]{D10}
D\'ecoppet T.D., On the dualizability of fusion 2-categoriess, \textit{Quantum Topol.}, {t}o appear, \href{https://arxiv.org/abs/2311.16827}{arXiv:2311.16827}.

\bibitem[10]{DY23}
D\'ecoppet T.D., Yu M., Fiber 2-functors and {T}ambara--{Y}amagami fusion
 2-categories, \href{https://arxiv.org/abs/2306.08117}{arXiv:2306.08117}.

\bibitem[11]{DR}
Douglas C.L., Reutter D.J., Fusion 2-categories and a state-sum invariant for
 4-manifolds, \href{https://arxiv.org/abs/1812.11933}{arXiv:1812.11933}.

\bibitem[16]{ENO2}
Etingof P., Nikshych D., Ostrik V., Fusion categories and homotopy theory,
 \href{https://doi.org/10.4171/QT/6}{\textit{Quantum Topol.}} \textbf{1} (2010), 209--273, \href{https://arxiv.org/abs/0909.3140}{arXiv:0909.3140}.

\bibitem[18]{GJF}
Gaiotto D., Johnson-Freyd T., Condensations in higher categories,
 \href{https://arxiv.org/abs/1905.09566}{arXiv:1905.09566}.

\bibitem[22]{JFR3}
Johnson-Freyd T., Reutter D.J., Super-duper vector spaces~{I}, 2023, available
 at
 \url{https://homepages.uni-regensburg.de/~lum63364/ConferenceFFT/Reutter.pdf}.

\bibitem[24]{JFY}
Johnson-Freyd T., Yu M., Fusion 2-categories with no line operators are
 grouplike, \href{https://doi.org/10.1017/S0004972721000095}{\textit{Bull. Aust. Math. Soc.}} \textbf{104} (2021), 434--442,
 \href{https://arxiv.org/abs/2010.07950}{arXiv:2010.07950}.

\bibitem[30]{L2}
Lurie J., Higher algebra, 2017, available at
 \url{https://www.math.ias.edu/~lurie/papers/HA.pdf}.

\bibitem[32]{Pstr}
Pstr\c{a}gowski P., On dualizable objects in monoidal bicategories,
 \textit{Theory Appl. Categ.} \textbf{38} (2022), 257--310,
 \href{https://arxiv.org/abs/1411.6691}{arXiv:1411.6691}.

\bibitem[34]{SP}
Schommer-Pries C.J., The classification of two-dimensional extended topological
 field theories, Ph.D.~Thesis, {U}niversity of California, Berkeley, 2009,
 \href{https://arxiv.org/abs/1112.1000}{arXiv:1112.1000}.

\end{thebibliography}
\end{document}
